\chapter*{ВВЕДЕНИЕ}
\addcontentsline{toc}{chapter}{ВВЕДЕНИЕ}

Компьютерная графика стала неотъемлемой частью современного мира и с каждым годом находит свое применения все в новых аспектах жизни. Она применяется в кинематографе, игровой индустрии, архитектура, медицине и других областях.

Алгоритм трассировка лучей в последние годы активно применяется в компьютерной графике. В связи с этим производители графических карт (GPU) начали добавлять в свои устройства ядра, специализирующиеся на трассировке лучей (RT-ядра). В стороне не остались разработчики графических API и библиотек, которые активно добавляют в свои продукты поддержку этой технологии. Известными примерами могут стать DirectX Ray Tracing (DXR), Vulkan RTX, CUDA/OptiX и другие. Многосторонняя поддержка этой технологии позволяет создавать все более фотореалистичные изображения с большим разрешением и лучшим качеством.

\textbf{Целью} данной работы является разработка ПО для визуализации трехмерных сцен, состояющих из геометрических примитивов, а также проведения исследования для тестирования производительности разработанного ПО.

Для достижения поставленной цели необходимо выполнить следующие \textbf{задачи}:
\begin{itemize}
    \item выбрать алгоритм генерации изображения, модель освещения и затенения;
    \item описать доступные для размещения на сцене модели и формализовать их;
    \item построить схемы алгоритмов;
    \item выбрать средства реализации (платформу, язык программирования, графическую библиотеку, GUI библиотеку);
    \item реализовать спроектированное ПО с помощью выбранных средств;
    \item провести исследование производительноси разработанного ПО.
\end{itemize}

\cite{nvidia}